\newpage
\section{Numerical Methods}

This section describes the numerical methods used in the photon kinetics
module in OSIRIS. It begins with how the photons are advanced (Time Integration),
how photon and electron densities are deposited onto the grid (Density
Deposition), how their gradients are calculated (Gradient Deposition), how
the grid quantities are filtered (Smoothing and Filtering), and conventions
for the units. This section assumes that you've read the Theoretical Background 
section, which contains the equations of motion that will be discretized here.
To see the  code that is actually implemented, see the OSIRIS Implementation 
section.

\subsection{Time Integration}

Here we will cover....

\subsubsection{Discretizing the Equations of Motion}

The ray equations are derived in the Theoretical Background section and are only
restated here. In OSIRIS units (lengths normalized by $c/\omega_{p0}$, times
normalized by $1/\omega_{p0}$, density normalized by $n_0$), a photon 
macro-particle with position $\mathbf{x}$ and wavenumber $\mathbf{k}$ obeys
\begin{equation*}
    \frac{d\mathbf{x}}{dt} = \frac{\mathbf{k}}{\omega}\,,
    \qquad
    \frac{d\mathbf{k}}{dt} = -\frac{1}{2\,\omega}\nabla n_e\,,
    \qquad
    \omega = \sqrt{|\mathbf{k}|^{2} + n_e(\mathbf{x})}\,.
    \label{eq:ray_equations}
\end{equation*}

\paragraph{What has to be discretized.}
Time is advanced in  ...

\paragraph{What the integrator should do.} We ultimately care about how this
code performs for the application of photon acceleration, in which the quantity
we care about is the wavenumber, $\mathbf{k}$. Any kind of numerical drift in
$\mathbf{k}$ will therefore appear as an unphysical frequency drift. The 
integrator should:
\begin{enumerate}

    \item \textbf{Not let $\omega$ drift.} In a stationary medium, $\omega$ is
    exactly conserved, ...

    \item \textbf{Be second order accurate.} We'd like to match the accuracy of
    the rest of the OSIRIS code, so this rules out integrators like RK4.

    \item \textbf{Use as little gathers as possible.} This is basically expensive

\end{enumerate}
The comparison across the next few sections evaluates how each integration 
scheme meets the above requirements.

\subsubsection{Naive Leapfrog Implementation}
We start from a Hamiltonian which is separated into a kinetic part that depends
only on momentum and a potential part that depends only on position: 
$H(\mathbf{q}, \mathbf{p}) = T(\mathbf{p}) + V(\mathbf{x})$.

In the above cases where the Hamiltonian is separable, the equations of motion
are given by
\begin{align*}
    \dot{\mathbf{q}} &= \nabla_{\mathbf{p}} \, H(\mathbf{q}, \mathbf{p}) = \nabla_
    {\mathbf{p}} \, T(\mathbf{p}), \\
    \dot{\mathbf{p}} &= -\nabla_{\mathbf{q}} \, H(\mathbf{q}, \mathbf{p}) = -\nabla_
    {\mathbf{q}} \, V(\mathbf{q}). 
\end{align*}
Each piece can be integrated exactly on its own: a solution for the coordinates
can be found if the momenta are kept constant and vice versa. These solutions 
can be combined in a time-symmetric fashion to give you the leapfrog method. There
are a few different types:
\begin{itemize}

    \item \textbf{Staggered Leapfrog.} This is the standard in PIC codes and what
    I'll be referring to from this point forward. This framing updates the
    position and velocity at interleaved times.

    \item \textbf{Velocity Verlet/Position Verlet.} These are mathematically
    equivalent to staggered Leapfrog, but synchronizes positions and velocities
    to the same integer timesteps. The difference between the velocity or
    position Verlet algorithms is simply the order of the updates.

    \item \textbf{Central-Difference Leapfrog.}

\end{itemize}

Leapfrog is the standard in PIC codes for a few reasons:
\begin{enumerate}

    \item \textbf{Time-reversible:} if you advance the solution from some
    initial condition to a future point, you can then make that point your new
    initial condition and step back to where you started.

    \item \textbf{Symplectic:} preserves the phase-space area/volumn (in 
    canonical units).

    \item \textbf{Explicit:} calculates a future state directly with the 
    current available information, at the cost of being conditionally stable
    and requiring smaller timesteps.

\end{enumerate}

\paragraph{Why leapfrog won't work for this case.} The Hamiltonian for a 
photon in this situation is $H = \omega = \sqrt{k^2 + n_e(x)}$, which isn't 
separable. $\omega$ depends on both $x$ and $k$, so neither $\dot{x} = k/\omega$
or $\dot{k} = -(1/2\omega)\nabla n_e$ can be written as a function of one 
variable alone. Because leapfrog isn't even well-defined or applicable in its
usual form this way, we look for other integrators.

You can see the consequences of this in the following single-photon test

\subsubsection{The Predictor-Corrector Method}
\paragraph{Why not a symplectic or implicit method?} A simple way to keep the
symplectic property for a non-separable Hamiltonian (while also maintaining only
second-order accuracy) is the lowest-order symplectic Runge-Kutta method. 
To turn RK2 symplectic, one can use ...., which as a side effect also makes the
integrator implicit (the implicit midpoint rule):

EQN HERE

\paragraph{The scheme.} The scheme chosen for the photon kinetic module is an 
explicit trapezoidal predictor-corrector method (Heun's method). In general, 
these update equations are given by:


\subsubsection{Accuracy and Stability}


\subsection{Density Deposition}


\subsection{Gradient Deposition}
\subsubsection{Old Method: Direct Gradient Gather}
\subsubsection{New Method: Deposit-then-Gather}
\subsubsection{Sign Convention and Equivalence}


\subsection{Smoothing and Filtering}
\subsection{Normalization and Units}

